% MSc thesis context (not a project changelog).
\subsection{Relation to the MSc thesis}
\label{S7_Appendix_arc}

The accompanying MSc thesis~\cite{elmgreen2026thesis} developed unsupervised mouse sleep staging with mixture priors on variational autoencoder latents.
That work used a \emph{single} MSSV laboratory (Copenhagen, laboratory~3): a curated ten-mouse subset, expert-engineered spectral features for classical HMM baselines, and \textbf{encoder+decoder} subject conditioning (embeddings in both encoder and decoder).
Population leave-one-subject-out prior NMI reached about 0.70 when staging from the cGMVAE mixture prior; substage discovery at higher mixture counts used the same in-laboratory protocol (thesis Figs~27--30).

This report changes three elements for cross-laboratory evaluation.
First, \textbf{decoder-only} conditioning: subject embeddings enter the decoder only so held-out mice can be staged from encoder means and the mixture prior without subject identifiers (\nameref{S3_Appendix}).
Second, a \textbf{quality-screened 20-mouse cohort} from laboratories 2, 3, and 5 with per-site spectral preprocessing locks (\nameref{S4_Appendix}) and four-fold leave-mice-out cross-validation (Tables~\ref{table1}--\ref{table2}).
Third, a \textbf{spectral CNN front-end} and end-to-end warm sticky HMM--GMM prior replace the thesis two-step pipeline (cGMVAE latents followed by MAR-HMM decoding on expert features).

Synthetic HMM validation and raw time-domain failure modes from the thesis are summarised in \nameref{S1_Appendix_synth} and \nameref{S2_Appendix}.
Planned extensions (laboratory conditioning, richer subject/lab disentanglement, and raw-signal CNN encoders) are outlined in \nameref{S13_Appendix_future}.
